
\chapter{概率与统计学}
 \label{chapter::basic-prob-append}
 
 
本书假定读者已具备概率论与统计推断方面的基础知识。因此，本章并非对概率与统计内容的全面综述。为便于查阅，我将回顾那些对正文理解至关重要的核心概念。  
 
 

\section{概率}



\subsection{皮尔逊相关系数与多重相关系数平方}
\label{sec::measure-pearson-r2}

对于两个随机变量 $Y$ 和 $X$，定义皮尔逊相关系数为
$$
\rho_{YX} = \frac{  \cov(Y, X)  }{  \sqrt{  \var(Y)   \var(X) }   }
$$
该系数用于衡量 $Y$ 对 $X$ 的线性依赖程度。该定义在 $Y$ 与 $X$ 之间是对称的，即
$$
\rho_{YX}  = \rho_{XY} .
$$


设随机变量 $Y$ 与随机向量 $X$，定义多重相关系数的平方为
$$
R^2_{YX} = \textup{corr}^2(Y, X )
= \frac{   \cov(Y, X) \cov(X)^{-1} \cov(X, Y)  }{  \var(Y) }
$$
其中 $\cov(Y, X) $ 是一个行向量，而 $\cov(X, Y)$ 是一个列向量。该值同样用于衡量 $Y$ 对 $X$ 的线性依赖程度。不过此定义在 $Y$ 与 $X$ 之间并不具备对称性。



\subsection{多元正态随机向量}
\label{sec::multivariate-normal}
 
一个多元正态随机向量 $\textsc{N}(\mu, \Sigma)$ 由其均值向量 $\mu$ 与协方差矩阵 $\Sigma$ 唯一确定。将其划分为两个部分：
\[
\left(\begin{array}{c}
Y_{1}\\
Y_{2}
\end{array}\right)\sim\textsc{N}\left(\left(\begin{array}{c}
\mu_{1}\\
\mu_{2}
\end{array}\right),\left(\begin{array}{cc}
\Sigma_{11} & \Sigma_{12}\\
\Sigma_{21} & \Sigma_{22}
\end{array}\right)\right).
\]

首先，各边际分布均呈正态分布：
$$
Y_{1}  \sim\textsc{N}\left(\mu_{1},\Sigma_{11}\right),\quad 
Y_{2}  \sim\textsc{N}\left(\mu_{2},\Sigma_{22}\right).
$$
其次，若 $\Sigma_{22}$ 为正定矩阵，则条件分布同样呈正态分布：
\[
Y_{1}\mid Y_{2}=y_{2}\sim\textsc{N}\left(\mu_{1}+\Sigma_{12}\Sigma_{22}^{-1}(y_{2}-\mu_{2}),\Sigma_{11}-\Sigma_{12}\Sigma_{22}^{-1}\Sigma_{21}\right).
\]
 

 
 
 \subsection{$\chi^2$分布与$t$分布}
\label{subsection::chi2-t-distribution}



假设 $X_1,\ldots, X_n$ 是独立同分布且服从 $\N01$ 分布的随机变量。那么
$$
\sumn X_i^2 
$$
服从自由度为 $n$ 的 $\chi^2_n$ 分布。该 $\chi^2_n$ 分布的均值为 $n$，方差为 $2n$。 



假设$X\sim \N01, Q_n\sim \chi^2_n$且$X\ind Q_n$。那么
$$
\frac{X}{  \sqrt{Q_n / n} }
$$
服从自由度为$n$的$t_n$分布。若$n>1$，则$t_n$分布的均值为0。当$n=1$时，$t_1$分布也被称为柯西分布，该分布没有有限的均值。 







\subsection{柯西-施瓦茨不等式}\label{sec::cauchy-schwarz-inequality}

 
柯西-施瓦茨不等式存在多种形式。对于两个随机变量 $A$ 和 $B$，我们有
$$
 | E(AB)  | \leq \sqrt{   E(A^2) E(B^2)  }
$$
当且仅当 $B = \beta A$（其中 $\beta$ 为某一常数）时等号成立。若将 $A$ 和 $B$ 的中心化以使其均值为0，则可得
$$
| \cov(A, B)  | \leq \sqrt{   \var(A) \var(B)  } 
$$
此时当且仅当 $B = \alpha + \beta A$（其中 $\alpha$ 与 $\beta$ 为常数）时等号成立。

当 $A$ 与 $B$ 分别为定义在有限集合 $\{ a_1, \ldots, a_n \}$ 与 $\{ b_1, \ldots, b_n \}$ 上的均匀随机变量时，我们有
$$
 \left|  \sumn a_i b_i  \right| \leq \sqrt{  \sumn a_i^2   \sumn b_i^2  }
$$
若存在一个常数 $\beta$ 使得对所有 $i$ 均有 $b_i =  \beta a_i$，则上述等式成立。 
 
 



\subsection{塔性质与方差分解}\label{appendix::tower-property}

给定随机变量或向量 $A, B, C$，我们有
$$
E(A) = E\{   E(A \mid B) \} 
$$
以及
$$
E(A\mid C) = E\{   E(A \mid B, C) \mid C \}. 
$$

 
给定一个随机变量$A$以及多个随机变量或向量$B, C$，我们有
$$
\var(A) = E\{   \var(A \mid B) \}  +  \var\{   E(A \mid B) \} 
$$
以及
$$
\var(A\mid C) = E\{   \var(A \mid B, C)  \mid C \}  +  \var\{   E(A \mid B, C)\mid C \}  .
$$
我在研究生阶段时，我的老师卡尔·莫里斯和乔·布利茨斯坦将这个公式称为{\it 夏娃定律}，原因在于该公式的形式恰好是“EVVE”。随后他们又将第一个塔性质称为{\it 亚当定律}。



同样，我们可以将协方差分解为
$$
\cov(A_1, A_2) = E\left\{ \cov(A_1, A_2\mid B)   \right\} + \cov\{ E(A_1\mid B), E(A_2\mid B)  \}
$$
以及
$$
\cov(A_1, A_2\mid C) = E\left\{ \cov(A_1, A_2\mid B, C)  \mid C \right\} + \cov\{ E(A_1\mid B, C), E(A_2\mid B, C) \mid C  \}.
$$


\subsection{极限定理}


\begin{definition}
[概率收敛]\label{def::converge-in-prob}
若对于任意 $\varepsilon > 0$，当 $n\rightarrow \infty$ 时均有
$$
\pr (   |X_n - X| > \varepsilon  ) \rightarrow 0
$$
成立，则称一列随机变量 $(X_n)_{n\geq 1}$ 依概率收敛于 $X$。 
\end{definition}



\begin{definition}
[分布收敛]\label{def::converge-in-distribution}
若对于 $ \pr(X\leq x) $ 的所有连续点 $x$，当 $n\rightarrow \infty$ 时均有$$
\pr(X_n \leq x) \rightarrow \pr(X\leq x)
$$成立，则称一列随机变量 $(X_n)_{n\geq 1}$ 依分布收敛于 $X$。 
\end{definition}


概率收敛的强度要高于分布收敛。定义\ref{def::converge-in-prob}与\ref{def::converge-in-distribution}有助于阐述关于独立同分布（IID）随机变量的样本平均值的以下两个基本定理。 



\begin{theorem}[大数定律]
\label{thm::lln}
若$X_1,\ldots, X_n \iidsim X$且满足$E| X| <\infty$，则$\bar{X} = n^{-1} \sumn X_i \rightarrow E(X)$依概率成立。 
\end{theorem}


定理\ref{thm::lln}中所描述的大数定律指出：在极限情况下，样本均值会趋近于总体均值。 

\begin{theorem}[中心极限定理（CLT）]
\label{thm::clt}
若 $X_1,\ldots, X_n \iidsim X$ 且满足 $\var (X) <\infty$，则 
$$ 
\frac{   \bar{X} - E(X)   }
{ \sqrt{ \var (X) /n    } } \rightarrow \N01
$$ 
在分布意义上成立。 
\end{theorem}

定理\ref{thm::clt}中的中心极限定理指出：在极限情况下，标准化后的样本均值接近于标准正态随机变量。 


定理\ref{thm::lln}与\ref{thm::clt}为简便起见，均假定随机变量满足独立同分布条件。  
此外，针对独立或弱相关随机变量\citep[e.g.,][]{durrett2019probability}的样本均值而言，也存在诸多大数定律与中心极限定理。 




\subsection{德尔塔法}\label{section::delta-method}


德尔塔方法是一种强有力的工具，可用于推导一个渐近服从正态分布的随机向量经非线性变换后仍呈渐近正态性。下文我将介绍德尔塔方法的一个特例。


\begin{theorem}
[德尔塔法]\label{thm::delta-method}
假设 $\sqrt{n} (X_n - \mu) \rightarrow \textsc{N}(0, \Sigma)$ 在分布意义上成立，且函数 $g(x)$ 在 $\mu$ 处的导数 $ \nabla g(\mu)   $ 不为零。那么
$$
\sqrt{n} \{  g(X_n) - g(\mu) \}  \rightarrow \textsc{N} (0,   (\nabla g(\mu) ) \tran  \Sigma \nabla g(\mu)  )
$$
在分布意义上也成立。
\end{theorem}  

我将略去定理\ref{thm::delta-method}的证明。根据一阶泰勒展开，该证明是显而易见的：
$$
g(X_n) - g(\mu) \approx  (\nabla g(\mu) ) \tran  (X_n - \mu) .
$$
因此，$\sqrt{n} \{  g(X_n) - g(\mu) \} $ 近似服从 $\textsc{N}(0, \Sigma)$ 的线性变换，即 $\textsc{N} (0,   (\nabla g(\mu) ) \tran  \Sigma \nabla g(\mu)  )$。 




作为示例，我们可以运用德尔塔方法来推导该比值与乘积的渐近正态性。 

\begin{example}
[该比值的渐近正态性]\label{example::asymptotic-normal-ratio}
假设\begin{equation}\label{eq::normal-ratio}
\sqrt{n} \begin{pmatrix}
Y_n - \mu_Y \\
X_n - \mu_X
\end{pmatrix}
\rightarrow 
\textsc{N}\left(
\begin{pmatrix}
0\\
0
\end{pmatrix},
\begin{pmatrix}
\sigma_Y^2 & \sigma_{YX} \\
 \sigma_{YX} &  \sigma_X^2 
\end{pmatrix}
\right)
\end{equation}依分布成立，且$\mu_X \neq 0.$。应用定理\ref{thm::delta-method}可得\begin{equation}\label{eq::ratio-normality}
\sqrt{n} \left(   \frac{Y_n}{X_n}  -  \frac{\mu_Y}{\mu_X} \right)
\rightarrow
\textsc{N}\left(0,    \frac{  \sigma_Y^2 }{ \mu_X^2 }  + \frac{ \mu_Y^2 \sigma_X^2}{ \mu_X^4}  - \frac{2\mu_Y \sigma_{YX}}{\mu_X^3} \right)
\end{equation}依分布成立。特别地，若$X_n$与$Y_n$渐近独立，即$\sigma_{YX} = 0$，则$Y_n/X_n$的渐近方差简化为$\sigma_Y^2 / \mu_X^2 + \mu_Y^2 \sigma_X^2 / \mu_X^4$。具体推导留作习题\ref{hw::asymptotic-normality-ratio}。
\end{example}


在示例 \ref{example::asymptotic-normal-ratio} 中，其渐近方差的计算过程稍显繁琐。一个更简便的记忆方式是基于下述近似公式：
\begin{equation}\label{eq::ratio-slutsky}
\frac{Y_n}{X_n}  -  \frac{\mu_Y}{\mu_X}  = \frac{Y_n -  \mu_Y/\mu_X \cdot X_n  }{X_n} \approx   \frac{Y_n -  \mu_Y/\mu_X \cdot X_n  }{\mu_X}, 
\end{equation}
因此，该比值的渐近方差就等于 $$
 \frac{Y_n -  \mu_Y/\mu_X \cdot X_n  }{\mu_X},
 $$ 的渐近方差，而后者是 $Y_n$ 与 $X_n$ 的线性组合。虽然斯卢茨基定理能够使 \eqref{eq::ratio-slutsky} 中的近似关系变得严谨，但该定理已超出了本书的讨论范围。 


\begin{example}
[乘积的渐近正态性]\label{example::asymptotic-normal-product}
假设 \eqref{eq::normal-ratio}。应用定理 \ref{thm::delta-method} 可得出 \begin{equation}\label{eq::product-normality}
\sqrt{n} \left(   X_n  Y_n  -  \mu_X \mu_Y \right)
\rightarrow
\textsc{N}\left(0,     \mu_Y^2 \sigma_X^2 + \mu_X^2 \sigma_Y^2 + 2 \mu_X \mu_Y \sigma_{XY} \right)
\end{equation} 在分布意义上成立。在 $X_n$ 与 $Y_n$ 渐近独立且 $\sigma_{YX} = 0$ 的特殊情况下，$X_n Y_n$ 的渐近方差可简化为 $\mu_Y^2 \sigma_X^2 + \mu_X^2 \sigma_Y^2$。具体细节则留作习题 \ref{hw::asymptotic-normality-product} 供读者自行推导。
\end{example}



\section{统计推断}
\label{sec::statistical-inference}

\subsection{点估计}


假设 $\theta$ 是我们感兴趣的参数。通常情况下，问题中还会出现一些我们不感兴趣的其他参数，用 $\eta$ 来表示。统计学家将 $\eta$ 称为{\it 干扰参数}。根据所获得的数据，我们可以计算出一个估计量 $\hat\theta$。在本书中，我们采用频率派的视角，即假定 $\theta$ 是一个固定的数值，而 $\hat\theta$ 则由于数据的随机性而具有随机性。一个估计量需要满足以下两个基本要求。 

\begin{definition}
[无偏性]\label{def::unbiasedness}
估计量 $\hat\theta$ 对于参数 $\theta$ 而言是{\it 无偏}的，当且仅当
$$
E(\hat\theta) = \theta
$$
对于 $\theta$ 与 $\eta$ 的所有可能取值均成立。 
\end{definition}



\begin{definition}
[一致性]\label{def::consistency}
估计量 $\hat\theta$ 对于参数 $\theta$ 而言是{\it 一致的}，当且仅当，样本量趋于无穷大时，对于所有的 $\theta$ 与 $\eta$ 的可能取值，都有$$
 \hat\theta  \rightarrow \theta
$$依概率收敛。 
\end{definition}


无偏性要求估计量的均值等于我们所关注的真实参数值。一致性则要求当样本量趋于无穷大时，估计量能够趋近于该真实参数值。无偏性并不必然导致一致性，反之，一致性也并不必然意味着无偏性。无偏性有时会成为一种限制条件，因为即便在一些较为简单的统计问题中，满足无偏性也可能难以实现。因此，在大多数统计问题中，一致性往往才是最基本的要求。 





\subsection{置信区间}



点估计量 $\hat\theta$ 是一个与真实参数 $\theta$ 有所差异的随机变量。统计学家通常希望找到这样一个区间：它能够以给定的概率范围包含真实参数。该区间是根据数据计算得出的，因此本身也是一个随机变量。 

\begin{definition}
[置信区间]\label{def::confidence-interval}
一个依赖于数据的区间 $[  \hat\theta_\textsc{l} , \hat\theta_\textsc{u} ] $ 若满足以下条件，则可被视为 $\theta$ 的置信区间，且其覆盖概率至少为 $1-\alpha$：
$$
\pr(    \hat\theta_\textsc{l}    \leq \theta \leq  \hat\theta_\textsc{u}  ) \geq 1-\alpha 
$$
即对于所有可能的 $\theta$ 与 $\eta$ 取值均成立。  
\end{definition}


\begin{definition}
[渐近置信区间]\label{def::asymptotic-confidence-interval}
一个依赖于数据的区间 $[  \hat\theta_\textsc{l} , \hat\theta_\textsc{u} ] $ 若满足以下条件，则可被视为 $\theta$ 的渐近置信区间，且其覆盖概率至少为 $1-\alpha$：
$$
\pr(    \hat\theta_\textsc{l}   \leq \theta \leq  \hat\theta_\textsc{u}  )   \rightarrow    1-\alpha'   ,\qquad \text{ as } n\rightarrow \infty
$$
此时 $\alpha '  \leq \alpha$，该条件对 $\theta$ 与 $\eta$ 的所有可能取值均成立。     
\end{definition}


一个常见的选择是 $\alpha = 0.05$。  
在定义 \ref{def::confidence-interval} 和 \ref{def::asymptotic-confidence-interval} 中，覆盖概率可以大于名义水平 $1-\alpha$。也就是说，这些定义允许覆盖度过高，但不允许覆盖度不足。当出现覆盖度过高的情况时，我们称该置信区间是保守的。当然，我们希望置信区间尽可能窄。否则，置信区间的定义便可以随意设定了。 



\subsection{假设检验}


  许多实际应用问题都可以表述为对某一假设进行检验：
  $$
  H_0: \theta = 0 .
  $$
决策规则 $\phi$ 是数据的二元函数：若我们拒绝零假设 $H_0$，则 $\phi = 1$；若未能拒绝 $H_0$，则 $\phi = 0$。该检验的第一类错误率是指在零假设成立的情况下我们仍将其拒绝的概率。下文我将对该定义进行详细说明。 


\begin{definition}[第一类错误率]
\label{def::size-test}
当$H_0$成立时，将检验$\phi$的第一类错误率定义为在所有可能的$\theta$与$\eta$取值下该概率的最大值。
$$
\pr(\phi = 1)
$$     
\end{definition}
  

一个常见的做法是确保第一类错误率低于 $\alpha = 0.05$。该检验的第二类错误率是指在零假设不成立时未拒绝该假设的概率。下文我将对该定义进行说明。 


\begin{definition}[第二类错误率]
\label{def::type2-error-test}
当$H_0$不成立时，将检验$\phi$的第二类错误率定义为在所有可能的$\theta$与$\eta$取值下该概率的最大值。$$
\pr(\phi = 0)
$$     
\end{definition}


在$H_0$条件下，由于第一类错误率得到了控制，我们希望在$H_0$不成立时，第二类错误率能够尽可能低。 






\subsection{Wald型置信区间与检验}\label{section::wald-type-ci-test}


许多统计学问题都具有如下结构。我们感兴趣的参数为 $\theta$。首先，我们找到一个一致估计量 $\hat\theta$，该估计量以概率收敛于 $\theta$；随后证明该估计量在渐近意义上服从均值为 $\theta$、方差为 $v$ 的正态分布，其中 $v$ 可能同时依赖于 $\theta$ 以及干扰参数 $\eta$。接着，我们根据解析公式或第 \ref{sec::bootstrap} 章中介绍的自助法来求得 该量 的一致估计量 $\hat v$。

$\hat v$ 的平方根被称为 {\it 标准误差}。

最终，我们可构造出 $\theta$ 的 Wald型置信区间如下：
$$
\hat{\theta} \pm z_{1-\alpha/2} \sqrt{ \hat{v} }
$$
此处 $z_{1-\alpha/2}$ 表示标准正态随机变量的 $1-\alpha/2$ 分位数。

该区间以约 $1-\alpha$ 的概率覆盖 $\theta$。若此区间不包含某个特定值 $c$（例如 $c = 0$），我们便拒绝零假设 $H_0(c): \theta = c$，这一检验过程即被称为 Wald检验。 

 
 
 \subsection{构建置信集与检验零假设之间的对偶性}
 \label{section::duality-ci-testing}

考虑针对标量参数 $\theta$ 的统计推断问题。统计学中的一个基本结论就是：构造 $\theta$ 的置信集与对关于 $\theta$ 的零假设进行检验是等价的。这一关系通常被称为“构造置信集”与“检验零假设”之间的对偶性。 

第\ref{section::wald-type-ci-test}节已探讨了基于Wald型置信区间与检验的对偶性。事实上，这种对偶性普遍成立。假设$\hat\Theta$是参数$\theta$的一个$(1-\alpha)$水平置信集：
$$
\pr(  \theta \in \hat\Theta  ) \geq  1-\alpha.
$$
此时，若数值$c$并不属于集合$\hat\Theta$，我们便可拒绝零假设$H_0(c): \theta = c$。该检验方法是有效的：当$\theta$确实等于$c$时，第一类错误率满足$\pr(  \theta \not  \in \hat\Theta ) \leq  \alpha$，这一数值恰好符合预期。
反之，若我们对一系列零假设$H_0(c): \theta = c$逐一进行检验，便可得到一系列与$c$相对应的$p$值，即$p(c)$。随后，所有在显著性水平$\alpha$下未被拒绝的$c$值便构成了参数$\theta$的一个置信集：
$$
 \hat\Theta = 
\{  c:  p(c) \geq \alpha \} = \{  c: \text{ fail to reject } H_0(c) \text{ at level } \alpha \} . 
$$
该集合确实是一个有效的置信集，因为
$$
\pr( \theta \in  \hat\Theta ) = 
\pr\{  \text{fail to reject } H_0(\theta) \text{ at level } \alpha \} \geq   1-\alpha.
$$


在此我使用“置信集”这一术语而非“置信区间”，因为通过反转检验得到的 $ \hat\Theta$ 未必会形成一个区间。  
关于对偶性的应用，可参见第 \ref{sec::fc-problem} 章与第 \ref{sec::frt-invert-interval} 章的内容。 



\section{基于2×2表的推断}



\subsection{费舍尔精确检验}
\label{sec::fisher-exact-test}

费舍尔提出了一种针对统计模型下 $H_0: p_1 = p_0$ 的精确检验方法：
$$
n_{11} \sim \text{Binomial}(n_1, p_1),\quad
n_{01} \sim \text{Binomial}(n_0, p_0),\quad 
n_{11} \ind n_{01}.
$$
下表对相关数据进行了总结。

\begin{center}
\begin{tabular}{cccc}
\hline 
 & 1 & 0 & 行求和 \\
 \hline 
样本1 & $n_{11}$ & $n_{10}$ & $n_1$ \\
样本0 & $n_{01}$ & $n_{00}$ & $n_0$ \\
列求和 & $n_{\cdot 1}$ & $n_{\cdot 0}$ & $n$ \\
\hline 
\end{tabular}
\end{center}





他认为，求和式 $n_{11}  + n_{01}  =  n_{\cdot 1}$ 以及 $n_{10}  + n_{00}  =  n_{\cdot 0}$ 所包含的信息对于反映 $p_1$ 与 $p_0$ 之间的差异而言十分有限；在给定这些求和值的前提下，$n_{11}$ 服从超几何分布，该分布并不依赖于零假设 $H_0$ 下的未知参数 $p_1=p_0$：
$$
\pr(n_{11} = k )  = \frac{  \binom{ n_{\cdot 1} }{  k }  \binom{  n -  n_{\cdot 1} }{  n_1 - k   }  }{  \binom{  n }{ n_1 }  }.
$$
在 \ri{R} 中，函数 \ri{fisher.test} 实现了该检验。 



\subsection{利用二乘二表格进行估计}\label{section::asymptotic-inference-2x2}


基于\ref{sec::fisher-exact-test}节中所介绍的模型，我们可以借助样本频率来估计参数$p_1$与$p_0$：
$$
\hat p_1 = \frac{n_{11}}{n_1} ,\quad \hat p_0 = \frac{n_{01}}{n_0}.
$$
因此，我们也能借助相应的样本类比值来估计风险差、对数风险比以及对数优势比
\begin{eqnarray*}
\RD &=& p_1 - p_0,\\ 
\log \RR &=&\log \frac{ p_1}{p_0},\\ 
\log \OR &=& \log \frac{ p_1/(1-p_1)}{p_0/(1-p_0)}
\end{eqnarray*}
\begin{eqnarray*}
\hat \RD &=& \hat  p_1 -\hat  p_0,\\ 
\log \hat \RR &=& \log\frac{ \hat p_1}{\hat p_0},\\ 
\log \hat \OR &=&\log \frac{ \hat p_1/(1-\hat p_1)}{\hat p_0/(1-\hat p_0)}  = \log\frac{n_{11}n_{00}}{n_{10}n_{01}}.
\end{eqnarray*}
根据渐近近似法（参见习题\ref{hw::variance-estimators-2x2}），上述各参数的估计方差值分别为
\begin{eqnarray*}
&&\frac{ \hat  p_1 (1-\hat  p_1) }{n_1} + \frac{ \hat  p_0 (1-\hat  p_0) }{n_0},\\
&& \frac{   1-\hat  p_1 }{n_1 \hat  p_1} + \frac{   1-\hat  p_0 }{n_0 \hat  p_0 },\\
&& \frac{ 1 }{n_1  \hat  p_1 (1-\hat  p_1)} + \frac{ 1 }{n_0 \hat  p_0 (1-\hat  p_0)},
\end{eqnarray*}
由于风险比与优势比的值始终为正数，上述对数变换能够带来更符合正态分布特性的近似结果。 




\section{统计学中的两个著名问题}



\subsection{贝伦斯-费舍尔问题}\label{sec::bf-problem}



考虑一个双样本问题：有 $n_1$ 个单位接受了处理，另有 $n_0$ 个单位作为对照。假设接受处理的单位所对应的结果值 $\{  Y_i: Z_i = 1 \}$ 独立同分布地取自正态分布 $\textsc{N} (\mu_1, \sigma_1^2)$，而对照组的相应结果值 $\{  Y_i: Z_i = 0 \}$ 则独立同分布地取自正态分布 $\textsc{N} (\mu_0, \sigma_0^2)$。我们的目标是检验零假设 $H_0: \mu_1 = \mu_0$。 



先从较简单的情形开始，即 $\sigma_1^2  = \sigma_0^2.$。依照第 \ref{ch::frt-cre} 章的内容，令 $$
\hat{\bar{Y}}(1) = n_1^{-1} \sum_{Z_i =1}  Y_i ,\quad
\hat{\bar{Y}}(0) = n_0^{-1} \sum_{Z_i =0}  Y_i
$$ 分别表示在处理组与对照组中观测结果的样本均值。一个经典的结论是 $$
t_\textup{equal} = 
\frac{  \hat{\bar{Y}}(1) - \hat{\bar{Y}}(0)    }{  \sqrt{\frac{n}{n_1 n_0(n-2)}  \left[  \sum_{Z_i=1}  \{Y_i - \hat{\bar{Y}}(1)\}^2
+ \sum_{Z_i=0}  \{Y_i - \hat{\bar{Y}}(0)\}^2
    \right]  }    }   \sim t_{n-2}.
$$基于 $t_\textup{equal} $，我们可以构造出一个用于检验 $H_0.$ 的检验方法。


现在考虑较难的情形，允许$\sigma_1^2  $与$ \sigma_0^2.$不同。此时$t_\textup{equal}$的分布不再是$t_{n-2}.$。分别估计两组方差，可定义$$
t_\textup{unequal} = \frac{    \hat{\bar{Y}}(1)  -   \hat{\bar{Y}}(0)    }{  \sqrt{  \frac{  \hat{S}^2(1)  }{n_1}  + \frac{\hat{S}^2(0)}{n_0}  }   },
$$其中$$
\hat{S}^2(1) = (n_1 - 1)^{-1} \sum_{Z_i =1} \{ Y_i  -\hat{\bar{Y}}(1)   \}^2,\quad
\hat{S}^2(0) = (n_0 - 1)^{-1} \sum_{Z_i =0} \{ Y_i  -\hat{\bar{Y}}(0)   \}^2 
$$分别为处理组和对照组的结果样本方差。遗憾的是，$t_\textup{unequal} $的精确分布依赖这些未知方差。在不假设方差相等时检验$H_0$，就是著名的Behrens–Fisher问题。若样本量$n_1$和$n_0$都很大，中心极限定理保证$t_\textup{unequal} $近似服从$\N01$，从而可构造$H_0.$的近似检验。由对偶性，$\mu_1 - \mu_0$的大样本Wald型置信区间为$$
 \hat{\bar{Y}}(1)  -   \hat{\bar{Y}}(0)   \pm z_{1-\alpha/2}    \sqrt{  \frac{  \hat{S}^2(1)  }{n_1}  + \frac{\hat{S}^2(0)}{n_0}  }     
$$其中$z_{1-\alpha/2}$为标准正态分布的$1-\alpha/2$分位数。


 



\subsection{菲勒尔-克里西问题}\label{sec::fc-problem}

考虑这样一个双样本问题：共有$n_1$个单位接受了处理，另有$n_0$个单位作为对照。假设接受处理的单位所对应的结果$\{  Y_i: Z_i = 1 \}$均独立同分布地取自$\textsc{N} (\mu_1, 1)$；而对照单位所对应的结果$\{  Y_i: Z_i = 0 \}$则独立同分布地取自$\textsc{N} (\mu_0, 1)$。我们的目标是估计$\gamma = \mu_1 / \mu_0$。此时可借助$\hat\gamma = \hat{\bar{Y}}(1)  / \hat{\bar{Y}}(0) $来估计$\gamma$。不过该点估计量的分布较为复杂，因此难以据此构建出针对$\gamma$的置信区间。 



Fieller置信区间的构建方式可通过对一系列零假设进行检验来实现：即 $H_0(c): \gamma = c$。在 $H_0(c)$ 成立的前提下，我们有
$$
\frac{  \hat{\bar{Y}}(1)  - c \hat{\bar{Y}}(0) }{  \sqrt{ 1/n_1 + c^2/n_0  } } \sim \N01 
$$
这一推导过程促成了如下置信区间的提出：
$$
\left\{
c: \Big | \frac{  \hat{\bar{Y}}(1)  - c \hat{\bar{Y}}(0) }{  \sqrt{ 1/n_1 + c^2/n_0  } } \Big | \leq z_{1-\alpha/2}
\right\}
$$
其中 $z_{1-\alpha/2}$ 表示标准正态随机变量的 $1-\alpha/2$ 分位数。


\section{统计学中的蒙特卡洛方法}\label{sec::monte-carlo-method}


蒙特卡洛方法是统计学中一种极为有力的工具。我将介绍该方法在近似计算期望值或平均值方面的基本应用，这对于理解第\ref{ch::frt-cre}章中所介绍的FRT概念而言至关重要。 


倘若我们的目标是计算 $\theta = E\{  g(Y)\}$，其中 $Y$ 为一个随机变量，我们只需从 $Y$ 的分布中抽取独立同分布样本 $Y_1,\ldots, Y_n$，进而得到 $\theta$ 的矩估计量：
$$
\hat\theta = n^{-1} \sumn g(Y_i).
$$

作为一种特殊情况，$Y$服从定义在$\{ y_1, \ldots, y_N \}$上的均匀分布，且$\theta = N^{-1} \sum_{i=1}^N g(y_i)$。此时我们可以从$\{ y_1, \ldots, y_N \}$中抽取$n$个独立同分布样本$\{ Y_1, \ldots, Y_n \}$，进而得到上文所定义的矩估计量$\hat\theta $。这种抽样方式被称为{\it 有放回抽样}，它与第\ref{chapter::SRS-appendix}章中介绍的{\it 不放回抽样}有所不同。 



\section{自助法}
\label{sec::bootstrap}


 
推导复杂估计量的方差公式往往非常繁琐。\citet{Efron_1979}提出了自助法，将其作为一种通用的方差估计工具。自助法存在诸多变体\citep{davison1997bootstrap}。而在本书中，我们仅需用到最基本的一种——非参数自助法，我们将其简称为“自助法”。 


考虑如下一般情形：$$
Y_1,\ldots, Y_n \iidsim Y,
$$ 其中 $Y_i$ 可以是一个通用的随机元素，代表第 $i$ 个单位所观测到的数据。估计量 $\hat\theta$ 是观测数据的函数：$\hat\theta = T(Y_1, \ldots, Y_n)$。当 $T$ 是一个复杂的函数时，要获得 $\hat\theta$ 的方差或渐近方差可能并非易事。 

$\hat\theta$的不确定性源于从真实分布中进行的独立同分布抽样，即对$Y_1,\ldots, Y_n$进行抽样。虽然真实分布未知，但当样本量$n$足够大时，其经验分布版本$$
\hat{F}_n(y) = n^{-1} \sumn I(Y_i \leq y),
$$能够很好地近似该真实分布。如果我们认可这一近似关系，便可通过抽样$$
(Y_1^*,\ldots, Y_n^* ) \iidsim \hat{F}_n(y).
$$的方式来模拟$\hat\theta$的值。由于$\hat{F}_n(y) $是一种离散分布，每个观测数据点上的概率质量均为$1/n$，因此模拟$\hat\theta$的过程可简化为以下步骤： 
\begin{enumerate}
\item
从集合$\{Y_1,\ldots, Y_n \}$中以有放回抽样的方式抽取样本$(Y_1^*,\ldots, Y_n^* ) $。
\item
计算 $\hat{\theta}^* = T(Y_1^*, \ldots, Y_n^*)$。
\item
重复上述两个步骤$B$次，即可获得自助法重复样本$\{  \hat{\theta}^*_1, \ldots, \hat{\theta}^*_B \}$。
\end{enumerate}




我们可借助自助法重复样本的方差来近似估计 $\hat\theta$ 的（渐近）方差：
$$
\hat{V}_\text{boot} = (B-1)^{-1} \sum_{b=1}^B  (\hat{\theta}^*_b - \bar{\theta}^*)^2,
$$
此处 $\bar{\theta}^* = B^{-1}\sum_{b=1}^B \hat{\theta}^*_b $。基于正态分布近似的自助法置信区间可表示为：
$$
\hat\theta \pm  z_{1-\alpha/2} \sqrt{\hat{V}_\text{boot} },
$$
其中 $z_{1-\alpha/2}$ 为标准正态分布变量的 $1-\alpha/2$ 分位数。



我使用以下这个简单示例来说明自助法的基本原理。在从分布$\textsc{N}(1,1)$中抽取$n=100$个独立同分布的样本$Y_i$后，样本均值应当接近$1$，其方差则为$1/100 = 0.01$。经过$500$次模拟实验后发现，传统的方差估计量以及采用$B=200$次重抽样的自助法方差估计量，它们的平均值都接近于$0.01$。 

\begin{lstlisting}
> ## sample size
> n  = 100
> ## number of Monte Carlo simulations
> MC = 500
> ## number of bootstrap replicates B
> n.boot = 200
> simulation = replicate(MC, {
+   Y = rnorm(n, 1, 1)
+   boot.mu.hat = replicate(n.boot, {
+     index = sample(1:n, n, replace = TRUE)
+     mean(Y[index])
+   })
+   c(mean(Y), var(Y)/n, var(boot.mu.hat))
+ })
> ## summarize the results
> apply(simulation, 1, mean)
[1] 0.997602961 0.010006303 0.009921895
\end{lstlisting}




\section{作业习题}

\paragraph{独立但不同分布的数据}
\label{hw::inid}

假设各$Y_i$相互独立，对$i=1,\ldots, n$，其均值为$\mu_i$、方差为$\sigma_i^2$。关注的参数为$
\mu = n^{-1} \sum_{i=1}^n \mu_i.
$。证明$
\hat\mu = n^{-1} \sum_{i=1}^n Y_i
$是$\mu$的无偏估计量，并求其方差。证明通常用于IID数据的方差估计量$$
\hat{v} = \{  n(n-1) \}^{-1} \sum_{i=1}^n (Y_i -  \hat{\mu})^2
$$是$\hat\mu$方差的保守估计量，即其期望不小于真实方差：$$
E(\hat{v}) - \var( \hat\mu ) = \{  n(n-1) \}^{-1} \sum_{i=1}^n (\mu_i - \mu)^2 \geq 0. 
$$

备注：考虑一个更简单的情形，即对所有 $i=1,\ldots, n$，均有 $\mu_i = \mu$ 且 $\sigma_i^2 = \sigma^2$。此时样本均值是 $\mu$ 的无偏估计量，其方差为 $\sigma^2/n$。此外，对于方差 $\sigma^2/n$ 而言，一个无偏估计量即为 $\hat \sigma^2/n = \hat v$，其中 $\hat\sigma^2 = (n-1)^{-1} \sum_{i=1}^n (Y_i -  \hat{\mu})^2$。

该问题阐述了一个更为一般性的结论，适用于观测值相互独立但分布并不相同的情形。这一结论对第 \ref{chapter::mpe} 章中讨论的配对实验所对应的定理 \ref{thm::varianceest-mp} 具有重要意义。 





\paragraph{比值的渐近正态性}\label{hw::asymptotic-normality-ratio}


证明\eqref{eq::ratio-normality}。


\paragraph{乘积的渐近正态性}\label{hw::asymptotic-normality-product}


证明 \eqref{eq::product-normality}。


\paragraph{两个独立正态分布的乘积}\label{hw-product-Normals}

假设 $X\sim \textsc{N}(\mu_X, \sigma_X^2), Y\sim \textsc{N}(\mu_Y, \sigma_Y^2)$ 且 $X\ind Y$。试证明：
$$
\var(XY) = \sigma_X^2 \sigma_Y^2 +\mu_X^2 \sigma_Y^2 + \mu_Y^2 \sigma_X^2 .
$$



备注：该问题给出了示例 \ref{example::asymptotic-normal-product} 的非渐近形式。 




\paragraph{ 2×2表格中的方差估计量}\label{hw::variance-estimators-2x2}

运用德尔塔法来证明\ref{section::asymptotic-inference-2x2}节中所提到的方差估计量。 




\paragraph{费舍尔加权法}\label{hw::fisher-weighting}


假设我们有 $p$ 个针对同一参数 $\theta$ 的独立且无偏估计量 $\hat{\theta}_1, \ldots, \hat{\theta}_p$：
$$
E(\hat{\theta}_j) = \theta,\quad (j=1,\ldots, p).
$$
每个估计量 $\hat{\theta}_j$ 的方差为 $v_j$ $(j=1,\ldots, p)$。 

构造一个新的估计量$$
\hat{\theta} = \sum_{j=1}^p w_j \hat{\theta}_j
$$该估计量包含一些非随机常数$w_j$。  
证明：若满足$\sum_{j=1}^p w_j = 1$，则$\hat{\theta} $是估计$\theta$的无偏估计量。进一步证明：在约束条件$\sum_{j=1}^p w_j = 1$下，能使$\hat{\theta} $的方差达到最小的最优常数$w_j$即为$$
w_j^* =  \frac{1/v_j}{ \sum_{j'=1}^p   1/v_{j'} }, \quad (j=1,\ldots, p).
$$  
最后证明：由此得到的估计量$\hat{\theta}^* = \sum_{j=1}^p w_j^* \hat{\theta}_j$的方差为$$
\var(\hat{\theta}^*) =  \frac{1 }{ \sum_{j=1}^p   1/v_{j} }. 
$$


备注：这被称为费舍尔加权法。最优权重与方差的倒数成正比。 


